
Le son est une onde mécanique, il se propage par compression-dilatation du
milieu matériel dans lequel il est émis. Visualiser, à l'aide d'un microphone
relié à un oscilloscope, le son émis par un instrument de musique puis par un
diapason.

\begin{figure}[H]
    \centering
    \begin{tikzpicture}[scale=0.8,every node/.style={font=\footnotesize}]
        \newdimen\y \y=1.5cm
        \draw[minorgridcolor,step=1mm] (0,-2) grid (8,2);
        \draw[majorgridcolor,step=1cm] (0,-2) grid (8,2);
        \draw[red,very thick]
            (0,0) sin (1,\y) cos (2,0) sin (3,-\y) cos (4,0) sin
            (5,\y) cos (6,0) sin (7,-\y) cos (8,0);
        \node[below] at (4,-2) {Onde émise par le diapason};
        \begin{scope}[xshift=8.6cm]
            \draw[minorgridcolor,step=1mm] (0,-2) grid (6,2);
            \draw[majorgridcolor,step=1cm] (0,-2) grid (6,2);
            \draw[blue,very thick,domain=0:18.85,samples=200]
                plot (0.318*\x,{sin(\x r)+sin(2*\x r)});
            \node[below] at (3,-2) {Onde émise par l'instrument};
        \end{scope}
        \begin{scope}[xshift=14.9cm]
            \node[right] {\includegraphics[scale=0.41]{act1_exp}};
        \end{scope}
    \end{tikzpicture}
\end{figure}
\begin{enumerate}
    \item Quel est le rôle du microphone~?
    \item Les ondes visualisées sont-elles périodiques~? Conclure.
    \item Comparer la forme des deux courbes observées.
    \item Déterminer la période de l'onde émise par le diapason, sachant que la
          sensibilité horizontale
          $S_{h}=\qty[per-mode=symbol]{0.5}{\ms\per\div}$.
    \item Calculer la fréquence $f$ de l'onde sonore émise par le diapason.
\end{enumerate}
